% Chapter IV — drop-in for the dissertation.
% Compile standalone:  pdflatex Chapter_IV_MORPH-OS.tex
% Or input from a book/report class after removing the standalone wrapper.
\documentclass[12pt,oneside]{book}

\usepackage[margin=1in]{geometry}
\usepackage{amsmath,amssymb,amsthm}
\usepackage{booktabs}
\usepackage{array}
\usepackage{listings}
\usepackage{xcolor}
\usepackage{hyperref}
\usepackage{microtype}
\usepackage{setspace}
\usepackage[T1]{fontenc}
\usepackage{lmodern}
\usepackage{url}

\onehalfspacing
\hypersetup{colorlinks=true,linkcolor=black,urlcolor=blue,citecolor=black}

\lstset{
  basicstyle=\ttfamily\footnotesize,
  breaklines=true,
  frame=single,
  backgroundcolor=\color{black!3},
  columns=fullflexible,
  keepspaces=true,
  showstringspaces=false
}

\providecommand{\R}{\mathbb{R}}
\providecommand{\T}{\mathbb{T}}
\newcommand{\abs}[1]{\lvert#1\rvert}
\newcommand{\ind}{\mathbf{1}}
\newcommand{\Cclass}{\mathcal{C}}
\newcommand{\Man}{\mathcal{M}}
\newcommand{\Ex}{\mathrm{Ex}}
\newcommand{\LEI}{\mathrm{LEI}}

\title{Chapter IV\\[0.4em]
\large Computational Operationalization of Extropic Semantic Codecs\\
in Reaction--Diffusion Topologies}
\author{Metta Thomas}
\date{2026}

\begin{document}
\frontmatter
\maketitle
\tableofcontents
\mainmatter

\chapter{Computational Operationalization of Extropic Semantic Codecs
in Reaction--Diffusion Topologies}

A \textbf{MORPH} (Minimal Ontic Reference, Preserved Homology) is an
ESC packet whose residual is sufficient to replay a Gray--Scott field
up to working-threshold $\beta_0$ on the periodic torus.

\section{Abstract and Methodological Scope}
\label{sec:scope}

This chapter operationalizes the Extropic Semantic Codec (ESC) on a
time-stepping Gray--Scott field. The theoretical apparatus of the
preceding chapters---extropy as novelty residual to a shared manifold,
contextual coherence, logical integrity---is here bound to a concrete
prior, a concrete packet, and a concrete homologue.

The substitution is declarative for imperative. An imperative codec
transmits the matrix that has already been computed. An ESC packet
transmits an ontology pointer, a manifold anchor, an extropic residual,
and a superlevel $H_0$ barcode. The decoder already possesses the
engine. Bit-for-bit parity of the spatial tensor is not required and is
not claimed.

Three quantities left schematic in the theoretical chapters are closed
in Section~\ref{sec:math}. Section~\ref{sec:homology} adds the
homologue: periodic superlevel $H_0$ of $V$, working threshold
$\tau=0.25$. Two artefacts are distinguished in Section~\ref{sec:impl}:
MORPH-OS is the specification; Listing~\ref{lst:l1} is a minimal L1
sketch. The laboratory protocol of Section~\ref{sec:emp} observes a
mature field and scores working-$\beta_0$ and $H_0$ bottleneck.

\paragraph{Taxonomy.}
Morphogen is the instrument. ESC is the packet format. A MORPH is the
homologue on the wire. MORPH-OS is the reference stack. MORPHOS is the
sibling IMAGE\_8 repository.

\section{Mathematical Formalization of the RD-Specific Codec}
\label{sec:math}

Encoder and decoder share engine \texttt{morphogen-gs-cpu-v2}, a Pearson
table $(F,k,D_U,D_V)$, and a canonical reset. Agreement is
\begin{equation}
  \label{eq:prior-hash}
  \texttt{prior\_hash}
  =\mathrm{sha256}(\texttt{engine}\mid\texttt{name}\mid F\mid k\mid D_U\mid D_V)[:16].
\end{equation}

\subsection{The shared manifold \texorpdfstring{\(\Man\)}{M}}
\label{sec:manifold}
\begin{equation}
  \label{eq:feature}
  x=\bigl(\overline{U},\,\overline{V},\,E,\,c_x,\,c_y,\,e,\,
  V_{10},\,V_{50},\,V_{90},\,H_{12}\bigr)\in\R^{21}.
\end{equation}
\begin{equation}
  \label{eq:manifold}
  \Man_{\Cclass}
  =\bigl\{
    x:\ \bigl\|(I-P_{T_p\Man})(x-\mu_{\Cclass})\bigr\|_2
    \le q_{0.95}(\Cclass)
  \bigr\}.
\end{equation}
Packet-declared class accuracy $1.00$; $z$-scored nearest neighbour $0.73$.

\subsection{Orthogonal novelty, LEI, coherence, integrity}
\label{sec:extropy}
\begin{equation}
  \label{eq:N}
  N(x)=\bigl\|(I-P_{T_p\Man})(x-p)\bigr\|_2.
\end{equation}
\begin{equation}
  \label{eq:lei}
  \LEI(x)=1-\cos\theta(x,\mu_{\Cclass}).
\end{equation}
\begin{equation}
  \label{eq:CL}
  \begin{aligned}
    C(x)&=\cos\bigl(x,\ \mathrm{mean}(\{x_i\}_{\Cclass})\bigr),\\
    L(x)&=\ind_{U,V\in[0,1]}\cdot
          \ind_{(F,k,D_U,D_V)\ \mathrm{legal}}.
  \end{aligned}
\end{equation}
\begin{equation}
  \label{eq:Ex}
  \Ex(x)=N(x)\cdot C(x)\cdot L(x).
\end{equation}

\subsection{Transformational delta \texorpdfstring{\(\Delta_u\)}{Delta u}}
\label{sec:delta}
\begin{equation}
  \label{eq:delta}
  \Delta_u(t)
  =\tfrac12\Bigl(
    W_1(H^U_t,H^U_{t-1})+W_1(H^V_t,H^V_{t-1})
  \Bigr).
\end{equation}

\subsection{Persistent homologous structures}
\label{sec:homology}
\begin{equation}
  \label{eq:beta0}
  \beta_0^{\mathrm{work}}(V)
  =\beta_0\bigl(\{V\ge 0.25\}\bigr)
  \quad\text{on }\T^2.
\end{equation}
Engine: \texttt{esc/homology.py}.

\subsection{Residual ladder}
\label{sec:ladder}
See the compiled chapter for Table~\ref{tab:ladder}.

\begin{table}[ht]
  \centering
  \caption{Residual ladder.}
  \label{tab:ladder}
  \begin{tabular}{@{}llp{0.42\textwidth}@{}}
    \toprule Mode & On the wire & Decoder \\
    \midrule
    L0 & pointer + anchor + \texttt{time\_step} & exact prior replay \\
    L1 & + disks in $[0,1]^2$ & seed disks, then integrate \\
    L2 & + $8\times 8$ $\overline{\Delta V}$ & upsample and add \\
    L3 & + $16\times 16$ chart & finer chart \\
    \bottomrule
  \end{tabular}
\end{table}

\section{Reference Implementation}
\label{sec:impl}

\begin{lstlisting}[caption={Minimal L1 sketch.},label={lst:l1}]
ENGINE = "morphogen-gs-cpu-v2"
prior_hash = sha256(engine|name|F|k|Du|Dv)[:16]
reset(); seed_disk(); step(); encode_L1(); decode_L1()
\end{lstlisting}

\section{Empirical Evaluation}
\label{sec:emp}

Full tables and prose are in the local compiled source
\texttt{Chapter\_IV\_MORPH-OS.tex} (artefact). Do not use this GitHub
copy as the sole source if it was truncated; prefer the artefact PDF.

\begin{thebibliography}{9}
\bibitem{morphos-repo}
M.~Thomas. MORPH-OS.
\url{https://github.com/IAmM3ta/MORPH-OS}, 2026.
\end{thebibliography}
\end{document}
